\documentclass[10pt,a4paper]{article}
\usepackage[margin=2cm]{geometry}
\usepackage{fancyvrb,booktabs,enumitem,amsmath,amssymb}
\fvset{fontsize=\scriptsize,frame=single}
\title{Lab Report -- Logical Reasoning for Planning}
\author{Apurva Verma}\date{}
\begin{document}\maketitle

\section*{Task 0: Planning problem specification}
(a) $I=\{\mathrm{At}(Robot,A),\mathrm{At}(Package,A)\}$. (b) $G=\{\mathrm{At}(Package,C)\}$.
(c,d) Actions:
\begin{center}\begin{tabular}{lll}\toprule Action & Preconditions & Effects\\\midrule
Move$(x,y)$, $(x,y)\in\{(A,B),(B,A),(B,C),(C,B)\}$ & At(Robot,$x$) & $\neg$At(Robot,$x$), At(Robot,$y$)\\
PickUp(Package,$l$) & At(Robot,$l$), At(Package,$l$) & $\neg$At(Package,$l$), Holding(Package)\\
Drop(Package,$l$) & At(Robot,$l$), Holding(Package) & $\neg$Holding(Package), At(Package,$l$)\\\bottomrule\end{tabular}\end{center}
\textbf{Initially applicable?} PickUp(Package,A): yes, both preconditions At(Robot,A) and At(Package,A) are in $I$. Drop(Package,C): no, At(Robot,C) and Holding(Package) are both false in $I$. (The planner confirms that the only initially applicable actions are Move(A,B) and PickUp(Package,A).) An action in the list is applicable only when $S\models \mathrm{Pre}(a)$.

\section*{Task 1: Manual plan}
Note: the lab's suggested sequence Move(A,B), PickUp(Package,B), \dots\ is not valid, since the package stays at A (PickUp at B needs At(Package,B)). Valid 4-step plan:
\begin{center}\begin{tabular}{ll}\toprule State & Facts\\\midrule
$S_0$ & At(Robot,A), At(Package,A)\\
$S_1$ = after PickUp(Package,A) & At(Robot,A), Holding(Package)\\
$S_2$ = after Move(A,B) & At(Robot,B), Holding(Package)\\
$S_3$ = after Move(B,C) & At(Robot,C), Holding(Package)\\
$S_4$ = after Drop(Package,C) & At(Robot,C), At(Package,C) $\models G$\\\bottomrule\end{tabular}\end{center}

\section*{Task 2: Prompt and generated program}
The prompt (the lab's suggested prompt plus the warehouse definition) is in the appendix and \texttt{prompts.txt}. \texttt{planner.py} was produced with LLM assistance from the prompts in the appendix and used without modification. Where the ideas appear: \emph{preconditions} $\to$ \texttt{applicable()} ($\texttt{pos\_pre}\subseteq S$ and \texttt{neg\_pre}$\cap S=\emptyset$); \emph{effects} $\to$ \texttt{apply\_action()} $=(S\setminus \texttt{neg\_eff})\cup\texttt{pos\_eff}$; \emph{goal} $\to$ \texttt{set(goal) <= s} when a state is dequeued, where planning stops; \emph{BFS} $\to$ the \texttt{deque} frontier with a \texttt{parent} dictionary, exploring all applicable actions level by level. Assumptions stated: closed-world assumption, only the package is movable, Move only along the listed edges.

\section*{Task 3: Tests}
\begin{center}\begin{tabular}{p{2.7cm}p{3.2cm}p{2cm}p{5cm}p{1.5cm}}\toprule
Test & Initial / goal & Plan found? & Plan & Valid?\\\midrule
A solvable & $I$ / At(Package,C) & yes & PickUp(Package,A), Move(A,B), Move(B,C), Drop(Package,C) & yes (re-executed, each precondition checked)\\
B no PickUp & $I$ / At(Package,C) & No plan found & -- & n/a\\
C1 extra Move(A,C), goal package at C & $I$ / At(Package,C) & yes & PickUp(Package,A), Move(A,C), Drop(Package,C) & yes\\
C2 extra Move(A,C), goal robot at C & $I$ / At(Robot,C) & yes & Move(A,C) & yes; package still at A\\
C3 extra Move(A,C), no PickUp & $I$ / At(Package,C) & No plan found & -- & n/a\\\bottomrule\end{tabular}\end{center}
Test B: the planner reports failure instead of inventing an action. Test C: when the robot alone can reach C, the package goal is still not satisfied (C3) and in C2 the robot reaches C without the package, so the planner does not conflate the two. An invalid plan (Move(A,B), Move(B,C), Drop) is rejected by the verifier (robot never held the package).

\section*{Task 4: Logic and Search}
Flow: Current state $\to$ check preconditions $\to$ \textbf{apply effects (remove negative, add positive)} $\to$ generate successor state $\to$ search over alternatives $\to$ goal? Logic (entailment $S\models \mathrm{Pre}(a)$ and the effect rules) decides which actions are possible and what the next state is; search (BFS with a visited set) decides in which order the applicable actions are tried and stops at the first state containing the goal. Logic determines what is possible; search determines what to try.

\section*{Task 5 (optional): LLM explanation versus independent verification}
The justification below was produced from prompt 3 in the appendix and compared with the transitions the program executed; they matched: PickUp needs At(Robot,A), At(Package,A) which hold in $S_0$; Move(A,B) needs At(Robot,A), true in $S_1$; Move(B,C) needs At(Robot,B), true in $S_2$; Drop needs At(Robot,C), Holding(Package), true in $S_3$. I trust (b) the independently executed state transitions more: they are computed mechanically from the action definitions, while an explanation is generated text that can sound right yet be wrong, for instance it could claim PickUp(Package,B) is valid as in the lab's example sequence. A generated explanation is not an independent verification.

\section*{Task 6--8: Prolog}
Run with SWI-Prolog 10.0.2 on \texttt{planner.pl}:
\begin{Verbatim}
?- can_move(a,b).   true
?- can_move(a,c).   false
?- valid_move(a,b). true
?- valid_move(b,c). true
?- valid_move(a,c). false
?- reduce_speed.    true
?- valid_path([a,b,c]). true
?- valid_path([a,c]).   false
\end{Verbatim}
\textbf{Task 6.} (a) \texttt{can\_move(a,b)} is true because the fact \texttt{connected(a,b)} exists and the rule \texttt{can\_move(X,Y) :- connected(X,Y)} derives it. (b) \texttt{can\_move(a,c)} is not established because there is no \texttt{connected(a,c)} fact and no rule that derives it (Prolog failure means ``not provable'' under the closed-world assumption, not ``proven false''). (c) The rule is the implication $\mathrm{Connected}(X,Y)\rightarrow\mathrm{CanMove}(X,Y)$ read right to left (head \texttt{:-} body), universally quantified over $X,Y$.

\textbf{Task 7.} \texttt{valid\_move(a,b)} and \texttt{valid\_move(b,c)} succeed, \texttt{valid\_move(a,c)} fails. \emph{Challenge:} the proposed Move(a,c) is not supported by the warehouse knowledge: no \texttt{connected(a,c)}, so Prolog rejects it, a check independent of the Python generator.

\textbf{Task 8.} \texttt{reduce\_speed} succeeds: \texttt{wet\_road} is a fact; \texttt{slippery} follows from it by the rule; \texttt{reduce\_speed} follows from \texttt{slippery}. Wet\_road $\Rightarrow$ (wet\_road $\rightarrow$ slippery) $\Rightarrow$ (slippery $\rightarrow$ reduce\_speed) $\Rightarrow$ reduce\_speed.

\section*{Think About It answers}
Applicability requires all preconditions to hold, not just that the action is listed (PickUp at B is listed but inapplicable). Logic determines what is possible; search determines what to try. A generated explanation is not an independent verification. Prolog is not identical to classical logic: it uses unification, backtracking, clause order and negation as failure.

\section*{Reflection Questions}
\begin{enumerate}[leftmargin=*]
\item Specifying preconditions and effects first fixes the meaning of ``correct'' so the generated code can be checked against it, and stops the LLM from inventing its own action semantics.
\item Without precondition checks the robot could Drop a package it is not holding, or PickUp at B when the package is at A, producing a ``plan'' that is impossible in the real world; in Test C the robot could wrongly ``deliver'' the package by just reaching C.
\item A plan can look reasonable (Move, Move, Drop) yet skip PickUp; validity means every action's preconditions hold in the state where it executes and the final state entails the goal.
\item The LLM wrote the state representation, \texttt{applicable}/\texttt{apply}, BFS and plan printing, plus the test harness and Prolog file, quickly.
\item What was verified independently: that the action definitions matched the specification, that re-executing the plan step by step succeeds, and the impossible, irrelevant-action and invalid-plan cases; it was also noticed that the lab's example sequence with PickUp at B is invalid.
\item In checking $S\models \mathrm{Pre}(a)$, applying effects, the goal test $S\models G$, and in the Prolog facts and rules (modus ponens chaining).
\item Planning is search in the space of states: states are sets of propositions, actions are the transitions (with logic deciding applicability), and BFS is the same algorithm used for path search in the previous module.
\end{enumerate}
\subsection*{Optional reflection (7.2)}
(1) A fact is an unconditional statement (\texttt{connected(a,b).}); a rule derives a head from body conditions (\texttt{can\_move(X,Y) :- connected(X,Y).}). (2) A query asks whether the goal is derivable from the facts and rules, i.e.\ whether it follows from the knowledge base. (3) Prolog checks the plan against an independent logical description, so a bug in the Python generator cannot validate itself. (4) An LLM-produced plan may be plausible but wrong; an independent verifier catches invalid actions without trusting the generator.

\section*{Appendix A: Prompts}\VerbatimInput{prompts.txt}
\section*{Appendix B: Output}
\subsection*{Python}\VerbatimInput{results_python.txt}
\subsection*{Prolog}\VerbatimInput{results_prolog.txt}
\section*{Appendix C: Code}
\subsection*{planner.py (LLM-generated, unmodified)}\VerbatimInput{planner.py}
\subsection*{planner.pl}\VerbatimInput{planner.pl}
\end{document}
